\section{Caso de negocio}\label{sec:negocio}

\subsection{Oportunidad: la IA ya planifica viajes, pero no inspira confianza}

La planificación de viajes con IA generativa ha pasado de curiosidad a costumbre en muy poco
tiempo. En Estados Unidos, el porcentaje de viajeros que usó IA para planificar, reservar o
moverse en destino subió del 33\,\% (primer semestre de 2025) al 56\,\% (febrero de
2026)~\cite{phocuswright2026aisurge}. El cuello de botella ya no es la adopción, sino la
\textbf{confianza}: solo el 8\,\% considera suficiente la respuesta de la IA y el 51\,\% acaba
abriendo las fuentes para comprobarla~\cite{phocuswright2026aisurge}. Además, uno de cada cuatro
viajeros ha recibido información desactualizada o incorrecta (25\,\% en el estudio de Amadeus, que
incluye España~\cite{amadeus2025connected}; 27\,\% según Global Rescue~\cite{globalrescue2026}). El caso extremo está
documentado: dos turistas recorrieron Perú en busca de un «Cañón Sagrado de Humantay» que un
\emph{chatbot} había inventado mezclando dos topónimos reales~\cite{aiid2025humantay}.

De ahí salen los tres problemas que aborda \kyrian{}:
\begin{enumerate}[label=\textbf{P\arabic*}]
  \item \textbf{Lugares inventados o desactualizados}: un modelo general responde con lo que
        recuerda, no con lo que puede citar.
  \item \textbf{Planes sin geografía}: listas de texto que no dicen si la mañana del día~2
        obliga a cruzar la ciudad dos veces.
  \item \textbf{Planes que no se pueden editar}: cambiar «la cena de esta noche por algo más
        barato» obliga a rehacer el itinerario.
\end{enumerate}

\subsection{Mercado}

El turismo internacional registró 1.520 millones de llegadas en 2025 (+4\,\%) y unos ingresos por
exportaciones de 2,2 billones de dólares~\cite{untourism2026barometer}. El mercado de viajes en
línea supera los 700.000 millones de dólares~\cite{statista2026onlinetravel}, y el segmento
específico de IA aplicada al turismo se estima en 4.300 millones de dólares en 2025, con un crecimiento anual
compuesto del 18,2\,\% hasta 17.300 millones en 2033~\cite{grandview2025aitourism}. Los usuarios
que planifican con IA, además, son el cliente de mayor valor: renta y gasto en viajes claramente
superiores a la media~\cite{phocuswright2026aisurge}.

\subsection{Competencia y diferenciación}

La Tabla~\ref{tab:competencia} resume los competidores más relevantes. Los grandes actores parten
de su inventario (hoteles, vuelos, reseñas) o de un modelo general; ninguno muestra, para cada
recomendación, la fuente y la licencia del dato que la respalda.

\begin{table}[t]
\centering\scriptsize
\caption{Competidores y posición de \kyrian{}.}\label{tab:competencia}
\begin{tabularx}{\columnwidth}{@{}L{1.75cm}L{2.2cm}Y@{}}
\toprule
\textbf{Producto} & \textbf{Fuente del plan} & \textbf{Hueco que deja} \\
\midrule
Mindtrip~\cite{mindtrip2025bvi} & Datos propietarios, mapa interactivo; versión para empresas (B2B) dirigida a oficinas de turismo & Corpus cerrado: no se puede auditar ni ceder \\
Google AI Mode (Canvas)~\cite{androidauthority2025canvas} & Datos de Maps, Flights y Hotels & Experimento en EE.\,UU.; asistente generalista \\
ChatGPT + apps de Expedia y Booking~\cite{lee2025chatgptapps} & Modelo general + inventario de la agencia en línea (OTA) & Orientado a vender alojamiento, no a conocer la ciudad \\
Tripadvisor AI~\cite{tripadvisor2023ai} & Reseñas de usuarios & Popularidad, no geografía; edición limitada \\
Layla, Wonderplan~\cite{fastio2026} & Modelo general & Borradores rápidos sin fundamentación \\
\midrule
\textbf{\kyrian{}} & \textbf{Corpus abierto por ciudad (Wikivoyage, Wikipedia, OpenStreetMap, Wikidata, Open-Meteo)} & \textbf{Cada tarjeta es un documento con identificador, coordenadas y licencia; el plan se dibuja y se edita sobre el mapa} \\
\bottomrule
\end{tabularx}
\end{table}

La propuesta de valor de \kyrian{} se apoya en cuatro elementos diferenciadores:
\begin{enumerate}[label=\textbf{D\arabic*}]
  \item \textbf{Fundamentación verificable}: el agente solo puede colocar en el itinerario
        lugares que la búsqueda le ha devuelto; un identificador que no viene del corpus se
        descarta antes de llegar al usuario (§\ref{sec:guardrails}).
  \item \textbf{Mapa como interfaz}: el itinerario se construye día a día sobre un mapa
        (MapLibre sobre OpenFreeMap), con barrios, hoteles y actividades geolocalizados.
  \item \textbf{Edición conversacional}: «cambia la cena» modifica una franja concreta del día mediante
        herramientas, sin rehacer el plan (§\ref{sec:agente}).
  \item \textbf{Coste marginal casi nulo}: la arquitectura 100\,\% \emph{serverless} no paga
        servidores en reposo y añadir una ciudad cuesta céntimos de cómputo (§\ref{sec:coste}).
\end{enumerate}

\subsection{Modelo de negocio y viabilidad}

\textbf{Clientes objetivo.} El usuario final es el viajero de escapada urbana (dos a cinco días en
una ciudad que no conoce), que hoy combina un asistente generalista, un mapa y varias guías y que
valora poder comprobar lo que se le recomienda. El cliente que paga, en la línea principal, es el
destino: oficinas de turismo, ayuntamientos y consorcios turísticos que quieren ofrecer en su web
un planificador con información fiable de su ciudad. Se contemplan dos líneas complementarias:
\begin{itemize}
  \item \textbf{B2B (venta a empresas) de marca blanca para destinos}. El mercado ya paga por ella: Mindtrip la vende a
        oficinas de turismo~\cite{mindtrip2025bvi}. Los textos y datos del corpus
        tienen licencia abierta y trazable, de modo que el destino puede auditar de dónde sale cada
        recomendación y enriquecer el corpus con su propia información; las fotos de alojamientos
        que proceden de la web de cada hotel se sustituirían por material del destino. Como
        hipótesis de precio, una licencia de unos 300--500\,€/mes por destino (alta, mantenimiento
        del corpus e inferencia incluidos) deja márgenes altos: el alta es un procedimiento
        automatizado (§\ref{sec:addcity}) y la inferencia es un coste variable. Con dos sesiones
        al mes por usuario (0,24\,\$), una licencia de 400\,€/mes cubre el gasto en IA de unos
        1.800 usuarios activos mensuales. La entrada al mercado sería un piloto gratuito de tres
        meses con una ciudad española de tamaño medio (su corpus se genera con \code{/add-city}
        en unas horas) para medir activación y conversión.
  \item \textbf{B2C (venta al consumidor) \emph{freemium} con afiliación}: el plan es gratuito y la monetización llega
        por comisión sobre actividades y alojamiento reservados desde las tarjetas; los programas
        de afiliación de actividades pagan en torno a un 8\,\%~\cite{track3602026} (cifra
        orientativa).
\end{itemize}

\textbf{Coste unitario.} Con el modelo de producción (Claude Haiku~4.5 en Bedrock, a 1\,\$ y
5\,\$ por millón de tokens de entrada y de salida, respectivamente), una llamada media del agente consume unos 3.360
tokens de entrada y 109 de salida, es decir, \textbf{unos 0,4 céntimos de dólar por llamada al
modelo}; un turno usa entre una y diez llamadas\footnote{Media
sobre las 295 llamadas de la evaluación de conversaciones de referencia (§\ref{sec:evalagente});
precios de la tabla versionada del repositorio.}. Un borrador completo de tres días cuesta 2--3
céntimos y cada edición alrededor de uno, de modo que una sesión típica
(borrador y una decena de ediciones) se queda en 10--18 céntimos de dólar, frente a unos 3\,€ de comisión
por una sola actividad de 40\,€ reservada. La cuota diaria por usuario (300.000 tokens por
defecto) acota el gasto a unos 0,35\,\$ por persona y día: el coste de IA crece con el uso, pero
nunca sin control. Los precios de AWS y de los modelos se dan en dólares, como se publican; el
coste de la infraestructura en reposo, en euros.

\begin{table}[h]
\centering\scriptsize
\caption{Escalado estimado del coste mensual (2 sesiones por usuario y mes a 0,12\,\$; infraestructura: solo orden de magnitud).}\label{tab:escalado}
\begin{tabularx}{\columnwidth}{@{}Yrrrr@{}}
\toprule
\textbf{Usuarios activos/mes} & \textbf{Sesiones} & \textbf{IA (\$)} & \textbf{Infra (\$)} & \textbf{Total (\$)} \\
\midrule
100 & 200 & 24 & $\approx$6 & $\approx$30 \\
1.000 & 2.000 & 240 & $\approx$10 & $\approx$250 \\
10.000 & 20.000 & 2.400 & $\approx$50 & $\approx$2.450 \\
\bottomrule
\end{tabularx}
\end{table}

La Tabla~\ref{tab:escalado} muestra que el coste es casi enteramente variable y proporcional al
uso: no hay escalones de capacidad que comprar. El umbral de rentabilidad en B2C es bajo:
\textbf{basta con que una de cada 25 sesiones acabe en la reserva de una actividad de 40\,€} para
cubrir el gasto en IA. La Tabla~\ref{tab:dafo} resume la posición estratégica del producto.

\begin{table}[h]
\centering\scriptsize
\caption{Análisis DAFO.}\label{tab:dafo}
\begin{tabularx}{\columnwidth}{@{}YY@{}}
\toprule
\textbf{Fortalezas} & \textbf{Debilidades} \\
\midrule
Fundamentación verificable por construcción; datos con licencia abierta y trazable; coste marginal mínimo; alta de ciudades automatizada; bilingüe. & Seis ciudades (mitigado: alta automatizada en horas); sin reserva integrada (prevista con afiliación); riqueza de OpenStreetMap y Wikivoyage desigual por ciudad (mitigado: puerta de calidad). \\
\midrule
\textbf{Oportunidades} & \textbf{Amenazas} \\
\midrule
Adopción de la IA en EE.\,UU. del 33 al 56\,\%; demanda de fuentes fiables; destinos que ya compran planificadores; distribución como aplicación dentro de los asistentes (protocolo MCP, \emph{Model Context Protocol}). & Grandes plataformas con datos propios (Google, Booking, Expedia); variación del precio de los modelos; mayores exigencias regulatorias de transparencia. \\
\bottomrule
\end{tabularx}
\end{table}

\textbf{Indicadores de negocio.} La Tabla~\ref{tab:kpinegocio} fija los indicadores con los que se
gestionaría el producto y su objetivo inicial (hipótesis para la fase de lanzamiento). Los que
dependen de viajes y turnos ya se miden con los eventos de negocio y las trazas
(§\ref{sec:metricas}), y los valores de los primeros 30 días con usuarios de prueba
(Tabla~\ref{tab:produccion}) ya cumplen los tres objetivos que se pueden medir.

\par\medskip\noindent\begin{minipage}{\columnwidth}
\centering\scriptsize
\captionof{table}{Indicadores de negocio, objetivos iniciales y valor medido en producción (30 días, 7 usuarios de prueba; --: se medirá tras el lanzamiento).}\label{tab:kpinegocio}
\begin{tabularx}{\columnwidth}{@{}L{1.75cm}Yrr@{}}
\toprule
\textbf{Indicador} & \textbf{Definición} & \textbf{Objetivo} & \textbf{Medido} \\
\midrule
Activación & Sesiones que generan un borrador / sesiones iniciadas & $\geq$40\,\% & -- \\
Viajes guardados & Viajes creados por usuario activo y mes & $\geq$1 & 2,9 \\
Implicación & Ediciones conversacionales por sesión & 3--10 & -- \\
Conversión & Clics de afiliación por tarjeta mostrada & $\geq$3\,\% & -- \\
Coste de IA & Gasto en modelos por usuario activo y mes & $\leq$0,30\,\$ & 0,29\,\$ \\
Calidad percibida & Turnos sin error ni reintento & $\geq$99\,\% & 99,7\,\%$^{*}$ \\
\bottomrule
\multicolumn{4}{@{}l@{}}{\tiny $^{*}$Sin las fichas pedidas con identificadores anteriores a una reindexación; con ellas, 97,1\,\% (§\ref{sec:produccion}).}
\end{tabularx}
\end{minipage}\par\medskip

En resumen, el mercado ya usa la IA para planificar; lo que falta es poder fiarse de ella, y ese
es el hueco que ocupa \kyrian{}.
